\section{Datasets}
\label{sec:datasets}

The datasets are assigned roles before interpretation. MetroPT and MetroPT2 are the primary event-level case studies because they contain ordered compressor telemetry and labelled failure periods \citep{Veloso2022MetroPT,MetroPT2Zenodo2023}. They support early-warning, alarm-episode, recurrence, and per-failure questions, but their independent event count is small. SCANIA Component X is an external vehicle-level stress test with irregular fleet histories, repair labels, and censoring \citep{Kharazian2025Scania}. Hydraulic Systems is a supplementary laboratory condition-monitoring demonstration with load-cycle component states \citep{Helwig2018HydraulicSystems}. SECOM is a supplementary high-dimensional production-quality demonstration with wafer-level yield labels \citep{McCann2008SECOM}.

This role assignment prevents a common mistake: row counts do not become independent failures. Table~\ref{tab:dataset-splits} is retained as a descriptive split inventory for the synthetic paper-asset cache. Table~\ref{tab:real-data-characteristics} records real-data source, independence unit, event or target source, and processed row counts. The independent unit column is the primary evidential quantity for empirical interpretation.

\input{../reports/paper/latex/dataset_split_summary.tex}

\begin{table}[t]
  \centering
  \caption{Registered real-data characteristics generated by the real-data matrix.}
  \label{tab:real-data-characteristics}
  \resizebox{\linewidth}{!}{\input{../artifacts/real_data_matrix/dataset_characteristics.tex}}
\end{table}

\begin{table}[t]
  \centering
  \caption{Dataset roles, sampling structures, independent units, verification status, and diagnostic result status.}
  \label{tab:evidence-scope}
  \resizebox{\linewidth}{!}{\input{../artifacts/real_data_matrix/evidence_scope.tex}}
\end{table}

The supported estimands differ. MetroPT-style rows address event-level early warning on temporal telemetry splits. SCANIA addresses vehicle-level repair-risk ranking on a published test split. Hydraulic Systems addresses cycle-level component-condition detection. SECOM addresses production-example yield-failure detection. Hydraulic Systems and SECOM are therefore not interpreted as event-level replications of the compressor task. This distinction is part of CLM-007: broader predictive claims require coherent independent units and comparable targets, not only successful data preparation.
